\documentclass[11pt,a4paper]{article}
\usepackage[utf8]{inputenc}
\usepackage[T1]{fontenc}
\usepackage[margin=0.7in, top=0.5in, bottom=0.5in]{geometry}
\usepackage{parskip}
\setlength{\parskip}{3pt}
\usepackage{lmodern}
\usepackage{tgheros}
\usepackage{xcolor}
\usepackage{titlesec}
\usepackage{enumitem}
\usepackage{hyperref}

\definecolor{slate}{HTML}{2c3e50}
\definecolor{teal}{HTML}{16a085}
\definecolor{rulecolor}{HTML}{bdc3c7}

\titleformat{\section}
  {\normalfont\large\bfseries\sffamily\color{slate}}
  {}{0em}{}[\color{rulecolor}\titlerule]

\hypersetup{
    colorlinks=true,
    linkcolor=teal,
    urlcolor=teal,
}

\begin{document}
\pagestyle{empty}

% --- Header ---
\begin{center}
    {\Large \bfseries \sffamily \color{slate} Research Statement} \\
    \vspace{0.2cm}
    {\small Yichao Jin, Ph.D. \quad | \quad \href{mailto:yichao.jin@utdallas.edu}{yichao.jin@utdallas.edu} \quad | \quad \href{https://yichao2022.github.io}{yichao2022.github.io}} \\
    {\small Methods: DCE, Structural Choice Modeling, LLM-Based Simulation, Behavioral Digital Twins}
\end{center}

\vspace{0.15cm}

\section*{Research Overview}

My research builds computational frameworks that translate behavioral models of patient and clinician decision-making into actionable, human-centered decision-support tools. I bridge behavioral economics, discrete choice methods, and generative AI to create systems that not only predict behavior but can be interacted with, audited, and deployed alongside real users.

\section*{The Behavioral Data Toolkit (BDT)}

My dissertation developed the \textbf{Behavioral Data Toolkit (BDT)}, a three-tier architecture that transforms Discrete Choice Experiment (DCE) data into high-fidelity synthetic agents:

\begin{itemize}[leftmargin=*, nosep]
    \item \textbf{Neural-Behavioral Ingestion} --- extracts latent preference parameters (risk tolerance, time discounting, trust sensitivity) from multi-modal choice data.
    \item \textbf{Probabilistic Decision Engine} --- a theory-agnostic core embedding structural econometric models (Random Utility / Random Regret Minimization) with Bayesian inference, enabling counterfactual simulation.
    \item \textbf{Synthetic Simulation Layer} --- deploys thousands of digital agents into virtual policy environments, generating synthetic micro-data for pre-screening interventions before field RCTs.
\end{itemize}

In a vaccination DCE (N=1,027, Wuhan, China), BDT reduced frontier-deviation MSE by 93.8\% compared to raw LLM agents and identified clinically meaningful heterogeneity in delay sensitivity by institutional trust --- a pattern conventional DCE analysis alone would not have revealed.

\section*{Alignment with Dr.~Raj's Group}

BDT is currently strongest at \emph{modeling} behavioral decision processes, but the question driving Dr.~Raj's group --- how to design decision-support tools that patients and clinicians actually trust and use --- is exactly where my work needs to go next. I see three concrete connections:

\begin{enumerate}[leftmargin=*, nosep]
    \item \textbf{Wearable health data as behavioral input.} BDT's ingestion layer currently handles experimental choice data; extending it to stream from passive sensor data (wearables, EHR logs) would allow preference models to update dynamically --- aligning with the group's interest in making continuous health data actionable for chronic condition management.
    \item \textbf{Red-teaming as validation.} The group's approach of stress-testing LLMs in clinical settings (e.g., red-teaming ChatGPT) provides a natural audit framework for BDT's synthetic agents. Behavioral digital twins must be robust to adversarial inputs before they can be trusted in shared decision-making.
    \item \textbf{Human-AI interaction design.} BDT currently outputs preference estimates and policy simulations, not interfaces. The group's expertise in designing interaction modalities that support metacognitive engagement --- helping users understand when and how to rely on AI recommendations --- would transform BDT from a modeling tool into a deployable decision-support system.
\end{enumerate}

\section*{Proposed Direction}

I propose to spend my postdoc adapting BDT's preference-elicitation framework into an interactive decision-support tool, validated through both behavioral experiments and clinical deployment studies --- contributing to the group's mission of making health data and algorithms genuinely actionable for the people who need them.

\vfill
\noindent{\small\sffamily\color{slate} Selected Publications:} {\small Jin, Y. (2026). \emph{Empirical-Frontier Regularization for LLM Synthetic Agents.} SSRN. \;|\; Jin, Y. (2026). \emph{Waiting Time as a Behavioral Barrier to Vaccination Uptake.} Job Market Paper.}

\end{document}
